\documentclass[aspectratio=169,10pt]{beamer}
\usetheme{default}
\setbeamertemplate{navigation symbols}{}
\usepackage{graphicx}\usepackage[T1]{fontenc}\usepackage{amsmath,amssymb}
\usepackage{listings}
\definecolor{ncsured}{HTML}{CC0000}\definecolor{ncsudk}{HTML}{990000}
\definecolor{ncsugray}{HTML}{6D6D6D}\definecolor{ncsuteal}{HTML}{008473}
\definecolor{codebg}{HTML}{F5F5F5}
\setbeamercolor{structure}{fg=ncsured}
\setbeamercolor{frametitle}{fg=white,bg=ncsured}
\setbeamercolor{title}{fg=ncsudk}
\setbeamercolor{itemize item}{fg=ncsured}
\setbeamercolor{itemize subitem}{fg=ncsured}
\setbeamerfont{frametitle}{series=\bfseries}
\setbeamertemplate{frametitle}{\nointerlineskip\vskip2pt%
  \begin{beamercolorbox}[wd=\paperwidth,ht=2.4ex,dp=1ex,leftskip=.3cm]{frametitle}%
  \small\bfseries\insertframetitle\end{beamercolorbox}}
\setbeamertemplate{footline}{\hbox{\begin{beamercolorbox}[wd=\paperwidth,ht=2.2ex,dp=1ex,leftskip=.3cm,rightskip=.3cm]{}%
  \tiny Optimal Market Making for Cryptocurrency Options \hfill Week 1 \hfill \insertframenumber\end{beamercolorbox}}}
\lstset{basicstyle=\ttfamily\scriptsize,backgroundcolor=\color{codebg},
  keywordstyle=\color{ncsured},commentstyle=\color{ncsugray},
  stringstyle=\color{ncsuteal},breaklines=true,columns=fullflexible,
  showstringspaces=false,frame=single,rulecolor=\color{ncsugray}}

\title{Week 1: Introduction \& Progress Review}
\subtitle{Roadmap, deliverables, grading, data sources, and resources}
\date{North Carolina State University \,\textbar\, Summer 2026}
\titlegraphic{
\includegraphics[width=4.2cm]{/Users/dendisuhubdy/Github/ncsu/assets/ncsu_logo.png}}
\begin{document}
\begin{frame}[plain]\titlepage\begin{center}\tiny Industry Mentor: Dendi Suhubdy (Bitwyre)  |  Guest: Fenni Kang (Coincall)\end{center}\end{frame}
\begin{frame}{Learning Objectives}
\begin{itemize}
  \item Understand the arc of the program and what you will build
  \item Be clear on the timeline, the three milestones, and how you are graded
  \item Know where the data comes from (Coincall) and how to pull it
  \item Have the Python + PyTorch and C++/pybind11 toolchains running
\end{itemize}
\end{frame}

\begin{frame}{Where We Are (Week 1 Progress Review)}
\begin{itemize}
  \item Maker vs. taker: we build the MAKER -- continuously quoting both sides
  \item The spine of the course is the Avellaneda-Stoikov model
  \item Week 1 status: environment setup, data access, and this orientation
  \item By Week 5 you ship a calibrated quoting engine (Milestone 1)
\end{itemize}
\end{frame}

\begin{frame}{Tooling \& How to Get Help}
\begin{itemize}
  \item Problem sets: Python + PyTorch (autograd does differentiation for you)
  \item Fast computation (Week 7): C++ exposed to Python via pybind11
  \item Mentorship is asynchronous (email/Telegram); tutors run day-to-day
  \item Office hours with tutors; come with a minimal reproducible snippet
\end{itemize}
\end{frame}

\begin{frame}{Timeline \& Deliverable Clarification}
\begin{itemize}
  \item 10 weeks. Week 1 introduction (done). Weeks 2-10 technical.
  \item Milestone 1 (Wk 5): Avellaneda-Stoikov engine, Python/PyTorch
  \item Milestone 2 (Wk 7): fast pricing/Greeks library, C++ via pybind11
  \item Milestone 3 (Wk 9): full options market-making backtest
  \item Final (Wk 10): one advanced extension + report (\textless{}=30pp) + 20-min talk
  \item Problem sets are due Sunday 23:59 ET; milestones Friday of the week
\end{itemize}
\end{frame}

\begin{frame}{Syllabus \& Grading Walkthrough}
\begin{itemize}
  \item Problem sets 30\%  |  Milestones (3) 30\%  |  Final 30\%  |  Participation 10\%
  \item Code graded on correctness, reproducibility, quality -- in that order
  \item Numerical-verification problems: right setup + right error magnitude
  \item Milestone 2 latency is a HARD GATE (sub-5ms surface) -- caps at 70 if missed
  \item Late PS: -20\% until Tue; nothing after. Milestones: no late submission.
  \item Set random seeds: reproducibility is explicitly graded
\end{itemize}
\end{frame}

\begin{frame}{Data Sources Discussion (Coincall)}
\begin{itemize}
  \item Exchange: Coincall.  REST base: https://api.coincall.com
  \item Provided as Parquet: BTC/ETH options L2 + trades, perp L2, funding, SVI
  \item Coincall orderbook is L2 (aggregated depth) -- there is NO L3 feed
  \item You can also pull live via the REST/WebSocket API (next slides)
  \item Auth: API key + HMAC-SHA256 signature; market data is largely public
  \item Always cross-check exact endpoints at https://docs.coincall.com/
\end{itemize}
\end{frame}

\begin{frame}{Additional Resources -- Math-Heavy Sections}
\begin{itemize}
  \item Stochastic control: Pham (2009), Ch. 3-4 (HJB, verification)
  \item Market making theory: Gueant (2016), Ch. 4; Cartea et al., Ch. 10-11
  \item Options/vol surface: Gatheral (2006); Gatheral-Jacquier (2014)
  \item Hedging: Taleb (1997), gamma scalping \& the gamma-theta identity
  \item El Aoud-Abergel (2015): multi-Greek option market making (Week 8)
  \item We DERIVE the key results in lecture and explain every modeling choice
\end{itemize}
\end{frame}

\begin{frame}{Illustration}
\begin{center}\includegraphics[width=0.86\textwidth]{/Users/dendisuhubdy/Github/ncsu/weekly_presentations/figures/week01.png}\end{center}
\end{frame}

\begin{frame}{Getting the Data from Coincall}
\begin{itemize}
  \item Base URL: https://api.coincall.com   |   Python SDK: github.com/CoincallExchange
  \item Signed request headers: X-CC-APIKEY, sign (HMAC-SHA256, upper hex), ts (ms), X-REQ-TS-DIFF
  \item Signing: sort params alphabetically, append uuid, ts, x-req-ts-diff, HMAC-SHA256 with secret
  \item Rate limits (per docs): orderbook 30/2s per user; public config 10/2s per IP
  \item Key endpoints you will use across the program:
  \begin{itemize}\item \texttt{\small Instruments:   GET /open/option/getInstruments/\{baseCurrency\}}\end{itemize}
  \begin{itemize}\item \texttt{\small Option chain:  GET /open/option/get/v1/\{index\}}\end{itemize}
  \begin{itemize}\item \texttt{\small Option book:   GET /open/option/order/orderbook/v1/\{symbol\}  (100 depth)}\end{itemize}
  \begin{itemize}\item \texttt{\small Option detail: GET /open/option/detail/v1/\{symbol\}  (mark, IV, delta/gamma/vega/theta)}\end{itemize}
  \begin{itemize}\item \texttt{\small Trades:        GET /open/option/trade/lasttrade/v1/\{symbol\}}\end{itemize}
  \begin{itemize}\item \texttt{\small Futures book:  GET /open/futures/order/orderbook/v1/\{symbol\}}\end{itemize}
  \begin{itemize}\item \texttt{\small Funding:       GET /open/public/fundingRate/v1?symbol=BTCUSD}\end{itemize}
  \begin{itemize}\item \texttt{\small Klines:        GET /open/option/market/kline/history/v1/\{optionName\}?period=h1}\end{itemize}
  \item WebSocket channels: /options/pricing, /options/orderbook, /options/lasttrade, /options/kline
  \item ALWAYS verify the exact path/params against https://docs.coincall.com/ -- the API evolves.
\end{itemize}
\vspace{2pt}\tiny\textit{Cross-check exact paths at https://docs.coincall.com/}
\end{frame}

\begin{frame}[fragile]{Code --- coincall\_client.py}
\begin{lstlisting}
import time, hmac, hashlib, uuid, requests

BASE = "https://api.coincall.com"

def signed_get(path, params=None, api_key="", api_secret=""):
    """Minimal Coincall signed GET. Public market-data endpoints often work
    unsigned; this shows the signed pattern used for account endpoints."""
    params = params or {}
    ts, window = str(int(time.time() * 1000)), "5000"
    items = sorted(params.items())                       # sort alphabetically
    payload = "&".join(f"{k}={v}" for k, v in items)
    payload += f"&uuid={uuid.uuid4().hex}&ts={ts}&x-req-ts-diff={window}"
    sign = hmac.new(api_secret.encode(), payload.encode(),
                    hashlib.sha256).hexdigest().upper()
    headers = {"X-CC-APIKEY": api_key, "sign": sign,
               "ts": ts, "X-REQ-TS-DIFF": window}
    r = requests.get(BASE + path, params=params, headers=headers, timeout=10)
    r.raise_for_status()
    return r.json()

# Example: list all BTC option instruments
print(signed_get("/open/option/getInstruments/BTC"))
\end{lstlisting}
\end{frame}
\begin{frame}{Glossary --- Key Terms}
\footnotesize\begin{description}
\item[\color{ncsured}Market maker] liquidity provider that continuously quotes bid and ask, earning the spread
\item[\color{ncsured}Taker] participant that crosses the spread for immediate execution
\item[\color{ncsured}Bid-ask spread] best ask minus best bid; the maker's revenue and risk premium
\item[\color{ncsured}Inventory] the maker's net position; the main source of risk
\item[\color{ncsured}Milestone] graded software deliverable, due in Weeks 5, 7, and 9
\item[\color{ncsured}PyTorch / autograd] Python ML library; automatic differentiation computes Greeks for you
\item[\color{ncsured}pybind11] library that exposes C++ functions to Python (used in Week 7)
\end{description}
\end{frame}

\begin{frame}{Reading}
\begin{itemize}
  \item Cartea-Jaikumar-Penalva, Ch. 1; Avellaneda-Stoikov (2008) -- skim
  \item Coincall API docs: https://docs.coincall.com/  (read 'Getting Started')
\end{itemize}
\end{frame}

\begin{frame}{References}
\scriptsize\begin{itemize}\setlength{\itemsep}{3pt}
  \item Avellaneda, M., \& Stoikov, S. (2008). High-frequency trading in a limit order book. Quantitative Finance, 8(3), 217-224.
  \item Gueant, O. (2016). The Financial Mathematics of Market Liquidity. Chapman \& Hall/CRC.
  \item Cartea, A., Jaikumar, S., \& Penalva, J. (2015). Algorithmic and High-Frequency Trading. Cambridge Univ. Press.
  \item Coincall API documentation. https://docs.coincall.com/
\end{itemize}
\end{frame}

\end{document}